\chapter{Fourier en tiempo discreto}\label{cap:fourier}

\enproceso{La DFT como cambio de base en $\C^N$ (de la tesis), con más cuidado: linealidad, desplazamiento, teorema de convolución, Parseval, periodicidad implícita, muestreo y aliasing, y la FFT para cualquier $N$ (no solo potencias de 2).}
